%% Verze pro jednostranný tisk:
% Okraje: levý 40mm, pravý 25mm, horní a dolní 25mm
% (ale pozor, LaTeX si sám přidává 1in)
\documentclass[12pt,a4paper,fleqn]{report}
\setlength\textwidth{145mm}
\setlength\textheight{247mm}
\setlength\oddsidemargin{15mm}
\setlength\evensidemargin{15mm}
\setlength\topmargin{0mm}
\setlength\headsep{0mm}
\setlength\headheight{0mm}
\let\openright=\clearpage

\usepackage[utf8]{inputenc}
\usepackage{graphicx}
\usepackage{amsthm}
% \usepackage{indentfirst}
\usepackage{amsfonts}
\usepackage{a4wide}
\usepackage[footnote]{acronym}
\usepackage{gnuplottex}
\usepackage{listings}
\renewcommand{\lstlistingname}{Algorithm}
\lstset{
	captionpos=b,
	breaklines=true,
	flexiblecolumns=false,
	literate={+}{{$+$}}1 {/}{{$/$}}1 {*}{{$*$}}1 {=}{{$=$}}1
             {>}{{$>$}}1 {<}{{$<$}}1 {\\}{{$\lambda$}}1
             {->}{{$\rightarrow$}}2 {>=}{{$\geq$}}2 {<-}{{$\leftarrow$}}2
             {<=}{{$\leq$}}2 {=>}{{$\Rightarrow$}}2 
             {|}{{$\mid$}}1
	}
% \lstloadlanguages{Haskell}
% \lstnewenvironment{code}
% 	{\lstset{}%
% 	  \csname lst@SetFirstLabel\endcsname}
% 	{\csname lst@SaveFirstLabel\endcsname}
% 	\lstset{
% 	  basicstyle=\small\ttfamily,
% 	  flexiblecolumns=false,
% 	  basewidth={0.5em,0.45em},
% 	  literate={+}{{$+$}}1 {/}{{$/$}}1 {*}{{$*$}}1 {=}{{$=$}}1
% 			   {>}{{$>$}}1 {<}{{$<$}}1 {\\}{{$\lambda$}}1
% 			   {\\\\}{{\char`\\\char`\\}}1
% 			   {->}{{$\rightarrow$}}2 {>=}{{$\geq$}}2 {<-}{{$\leftarrow$}}2
% 			   {<=}{{$\leq$}}2 {=>}{{$\Rightarrow$}}2 
% 			   {\ .}{{$\circ$}}2 {\ .\ }{{$\circ$}}2
% 			   {>>}{{>>}}2 {>>=}{{>>=}}2
% 			   {|}{{$\mid$}}1               
% 	}
	% more on http://www.haskell.org/haskellwiki/Literate_programming


\usepackage[unicode]{hyperref}
\hypersetup{pdftitle=Arimaa challenge - comparission study of MCTS versus alpha-beta methods}
\hypersetup{pdfauthor=Tomáš Jakl}


%%% Drobné úpravy stylu

% Tato makra přesvědčují mírně ošklivým trikem LaTeX, aby hlavičky kapitol
% sázel příčetněji a nevynechával nad nimi spoustu místa. Směle ignorujte.
\makeatletter
\def\@makechapterhead#1{
  {\parindent \z@ \raggedright \normalfont
   \Huge\bfseries \thechapter. #1
   \par\nobreak
   \vskip 20\p@
}}
\def\@makeschapterhead#1{
  {\parindent \z@ \raggedright \normalfont
   \Huge\bfseries #1
   \par\nobreak
   \vskip 20\p@
}}
\makeatother

% Toto makro definuje kapitolu, která není očíslovaná, ale je uvedena v obsahu.
\def\chapwithtoc#1{
\chapter*{#1}
\addcontentsline{toc}{chapter}{#1}
}

\begin{document}

% Trochu volnější nastavení dělení slov, než je default.
\lefthyphenmin=2
\righthyphenmin=2

%%% Titulní strana práce

\pagestyle{empty}
\begin{center}

\large

Charles University in Prague

\medskip

Faculty of Mathematics and Physics

\vfill

{\bf\Large BACHELOR THESIS}

\vfill

\centerline{\mbox{\includegraphics[width=60mm]{logo.eps}}}

\vfill
\vspace{5mm}

{\LARGE Tomáš Jakl}

\vspace{15mm}

% Název práce přesně podle zadání
{\LARGE\bfseries Arimaa challenge -- comparission study of MCTS versus alpha-beta methods}

\vfill

Department of Theoretical Computer Science and Mathematical Logic

\vfill

\begin{tabular}{rl}

Supervisor of the bachelor thesis: & Mgr. Vladan Majerech, Dr. \\
\noalign{\vspace{2mm}}
Study programme: & Computer Science \\
\noalign{\vspace{2mm}}
Specialization: & General Computer Science \\
\end{tabular}

\vfill

Prague 2011

\end{center}

\newpage

%%% TODO Na tomto místě mohou být napsána případná poděkování (vedoucímu práce,
%%% konzultantovi, tomu, kdo zapůjčil software, literaturu apod.)

\openright

\noindent
Poděkování
%I would like to thank to dr. Majerech Vladan for revealing this topic, which
%turned out to be more interesting than it seemed at first time.


%%% Strana s čestným prohlášením k bakalářské práci

\vglue 0pt plus 1fill

\noindent
I~declare that I~carried out this bachelor thesis independently, and only with
the cited sources, literature and other professional sources.

\medskip\noindent
I~understand that my work relates to the rights and obligations under the Act
No. 121/2000 Coll., the Copyright Act, as amended, in particular the fact that
the Charles University in Prague has the right to conclude a license agreement
on the use of this work as a school work pursuant to Section 60 paragraph 1 of
the Copyright Act.

\vspace{10mm}

\hbox{\hbox to 0.5\hsize{%
In Prague date \today
\hss}\hbox to 0.5\hsize{%
Tomáš Jakl
\hss}}

\vspace{20mm}
\newpage

%%% Povinná informační strana bakalářské práce (TODO)

\vbox to 0.5\vsize{
\setlength\parindent{0mm}
\setlength\parskip{5mm}

Title:
Arimaa challenge -- comparission study of MCTS versus alpha-beta methods

Author:
Tomáš Jakl

Department:
Department of Theoretical Computer Science and Mathematical Logic

Supervisor:
Mgr. Vladan Majerech, Dr., Department of Theoretical Computer Science and Mathematical Logic

Abstract:
In the world of chess programming the most successful algorithm for game tree
search is considered AlphaBeta search, however in game of Go it is Monte
Carlo Tree Search. Arimaa has similarities with both Go and Chess, but there
has been no successful program using Monte Carlo Tree Search so far. The main
goal of this thesis is to compare capabilites given by Monte Carlo Tree Search
algorithm and AlphaBeta search, both having the same evaluation function, in
game of Arimaa.

Keywords: Arimaa, Monte Carlo Tree Search, Alpha Beta search

\vss}\nobreak\vbox to 0.49\vsize{
\setlength\parindent{0mm}
\setlength\parskip{5mm}

Název práce:
Arimaa challenge -- srovnávací studie metod MCTS a alfa-beta

Autor:
Tomáš Jakl

Katedra:
Katedra teoretické informatiky a matematické logiky

Vedoucí bakalářské práce:
Mgr. Vladan Majerech, Dr., Katedra teoretické informatiky a matematické logiky

Abstrakt:
% TODO abstrakt v rozsahu 80-200 slov; nejedná se však o opis zadání bakalářské práce

Klíčová slova: Arimaa, Monte Carlo Tree Search, Alpha Beta search, UCT


\vss}

\newpage

\openright
\pagestyle{plain}
\setcounter{page}{1}
\tableofcontents

%%% Chapters
% \chapter*{Introduction}
% \addcontentsline{toc}{chapter}{Introduction}

\chapter{Introduction}
After Gasparov was defeated, new challenges has come. Defeat men in Go. Until
now, there were no significant success on standart 13$\times$13 board. But many
useful new way of playing games was invented. In ??? Monte Carlo methods were
successfully used to defeat all other computer players of Go. After that day
all successful programs were using Monte Carlo methods.
% TODO so many ???s

Tomáš Kozelek has shown in his work, that building Arimaa playing program using
MCTS is possible. In this work we will focus on comparing capabilities and
perspectives to the future given by AlphaBeta search and MCTS to game of
Arimaa.

After first computers were created it was always in human target to fight
with/compete human mind in every occasion. First significant result was shown
in 1997. IBM constructed/built computer with just one purpose, to defeat
the best human player in the game of chess. It took few years of development
and ??? milions of dolars to build computer. Fist time they wasn't successful
[???], but after some time and more effort they defeated Gasparov with computer
named Deep blue.
% TODO preformulovat a upresnit/rozsirit
[1] and [???]

We will introduce the game of Arimaa and describe two algorithms \ac{MCTS} and AlphaBeta search ... TODO


\section{Terminology}
\begin{description}
\item[Game bot] is game playing program.
\item[Game tree] for arbitrary game is tree with starting positions as root and
   with children of nodes as all possible consequent positions.
\item[Evaluation function] is function which estimates value of given position of the game. It can be used for example to compare which of two given positions is better.
\item[Minimax tree] for two player game is game tree limited to some depth with
added values in all nodes. Values are defined recursively. In leaf of the tree
is value calculated by evaluation function. In nonleaf node is the value
defined as the best value from nodes children from nodes active player point
of view.
\end{description}

\section{The Game of Arimaa}
The game of Arimaa is pretty new game. It is carefully designed to be hard to
play for computers, but easy to play for humans. Creator of the game Omar
Syed [1]

The game was carefully designed in order not to be possible to use methods
well known from Chess as game-ending tables or opening tables. (TODO
REWRITE:) Also to be significantly harder to precompute moves for huge
number moves to future and to efficiently decide which of two given position
is better.

(TODO REWRITE:) Omar says that gasparov was not oversmarted but overcomputed ...
In Arimaa it is significantly harder to precompute moves ??because?? and to
efficiently decide which of two given position is better ??because??.

Why we cannot use standart methods widely used in chess? (Section 2 in Kozeleks thesis)

\section{Rules of the game\cite{arimaa.com}}
Arimaa is two-player zero-sum game with perfect information. It is designed to
be possible to play it using the board and pieces from chess set. The starting
player has Gold color, second is Silver. In the zeroth turn Gold and then
Silver player each place all their pieces into first two (for Gold) or last two
(for Silver) lines of the board any way they consider appropriate. Piece set
consist of eight Rabbits, two Cats, two Dogs, two Horses, Camel and Elephant in
order from weakest to strongest.

%% TODO example picture of the starting position

Players are taking turns starting with Gold. In each turn player makes move,
which consist of one to four steps. If less than four steps were made it is
said that player passed. After move ends, position of the board must be
different from position before turn started.

Figure is considered frozen if on one of its neighbouring squares is opponents
stronger piece and on no figure sharing color with it. Adjacent or neighbouring
squares are those squares lying in one position to left, right, front or
backwards.

%% TODO example picture of the frozen piece and possibility to push and pull

To make step player chooses one of its non frozen figures and move it to one of
the free adjacent squares, with one exception -- rabbits cannot step backwards.
Instead of making simple step player can provide pull or push. Pulling is
almost identical to simple step, but after step is done player choose one of
the weaker piece which stayed on neighbouring of the square from which figure
moved, and move it to position from which our piece stepped. Pushing is a kind
of opposite to pulling. At first neighbouring weaker piece is moved to its free
adjacent square and then our piece is placed to weaker pieces former place.
Push and pull counts as 2 steps and can not be combined together.

On the board are four special squares called traps in positions \texttt{c3},
\texttt{c6}, \texttt{f3} and \texttt{f6}. After each step ends if any piece is
in trap and has no adjacent piece sharing color with him, it is trapped and
therefore removed from the board.

At end of the move is checked if game ended and it happens if one of the
following conditions is fulfilled:
\begin{enumerate}
\item one of the players rabbit has reached opponents front line (goalline), we
call it goal,
\item one of the players has no rabbit left, we call it elimination,
\item next active player cannot move, we call it immobilisation,
\item or the same position is repeated in third time in the row.
\end{enumerate}

Player wins if he scores goal or opponent is eliminated, is immobilised or
repeated position for third time.

%TODO figures with examples of pushing/pulling, immobilisations, trapping, ...

\section{Comparison to Go and Chess}
Because Arimaa is played with full chess set it is very natural to ask about
similarities with the game of Chess. As in Chess is in Arimaa also so easy to
ruin good position with just one bad move. For example stepping out of trap an
therefore let another piece to be trapped or unfreezing rabbit near goalline.
Unlike in Chess starting position is not predefined and there are about
$4.207\times10^{15}$ different possible openings which makes it hard for arimaa
bot programmer to use any kind of opening tables\cite{COX}.

In Arimaa it is also very hard to build good static evaluation function. Even
human players often do not know for example if in certain position it is better
to sacrifice Camel for trapping opponents Horse and Cat and often it depends on
many other factors. Building good evaluation function is very hard in game of
Go also, because it requires a lot of local searching to find territories and
determine their potential owners. On the other hand destroying good position by
making wrong step is in Go a lot harder.

In game of Go if we omit filling eye any random playing sequence leads to end,
which is not so easily achievable in Arimaa and Chess. Christ Cox showed in his
work that in Arimaa branching factor of average position is around 20,000. With
comparison in Chess it is only 35 and in Go 200 \cite{COX}.

% TODO? position in arimaa vs tactic in chess \cite{COX}

\section{Challenge}
Omar Syed decided to left a prize 10,000 USD for programmer or group of
programmers who develop Arimaa playing program which win Arimaa Computer
Championship and then defeat three chosen top human players before the year
2020. Omar Syed believes that this motivation will help further improvements in
area of AI game programming\cite{syed}.

So far computers were not even close to defeat one of the chosen human
players\cite{arimaa.com}. ... (TODO add facts)

\section{Object of research}
The main part of this work is to develop well documented Arimaa playing program
and try to answer the following questions:

\begin{enumerate}
\item Is MCTS competitive alpha beta search at all?
\item Is MCTS more promising engine than AlphaBeta search in the future with
      increasing number of cpus?
\item How important is evaluation function in MCTS compared to AlphaBeta?
\end{enumerate}
% TODO: which area should be compared? who is more perspective in future?


\chapter{Algorithms}
In this chapter we introduce two algorithms for searching minimax trees --
AlphaBeta search and Monte Carlo Tree Search (MCTS). Both algorithms became
successful and dominating in some field, AlphaBeta search in game of Chess and
MCTS in game of Go.

The basic algorithm to search game tree is Minimax search. It is searching
minimax tree in depth first search order and therefore the search do not ends
until all nodes at certain depth are explored. It is quite noneffective with
huge branching factor which Arimaa, Chess or Go have.

\section{Description of the AlphaBeta search}
AlphaBeta search has grown in popularity in game of Chess. Until now all
successful arimaa bots were using it with various enchants. AlphaBeta search is
natural optimisation of Minimax search. The main purpose of algorithm is to
reduce number of branches and nodes to be visited.

During the search we are trimming bounds (window) of the best minimax value.
The pseudocode of the algorithm is shown in~\ref{alg:alphabeta}.

\lstset{language=Python, caption=Pseudocode of the AlphaBeta search, label=alg:alphabeta}
\begin{lstlisting}
alphabeta (node, depth, alpha, beta):
    if  depth = 0 or is_terminal(node):
        return evaluate(node)
    if  is_maximizing(node):
        for each child of node
            score = alphabeta(child, depth-1, alpha, beta)
            alpha = max(alpha, score)

            if beta <= alpha: break  # Beta cut-off
        return alpha
    else:
        for each child of node:
            score = alphabeta(child, depth-1, alpha, beta)
            beta = min(beta, score)

            if beta <= alpha: break  # Alpha cut-off
        return beta
\end{lstlisting}

In maximizing nodes we are improving lower estimate (alpha) of the minimax
value and in minimizing nodes the upper estimate (beta). If a value from a
child forces these bounds to meet we know that better a approximation of this
node is not possible and so we return our estimate, we say we pruned/cut-off
the search in node.

It can be proven that if AlphaBeta finds solution for depth $n$, then it is the
best solution in MiniMax tree for depth $n$~\cite{knuth:alphabeta}. In optimal
case instead of examining $\mathcal O(b^d)$ nodes is examined only $\mathcal
O(b^{d/2})$ nodes, where $b$ is branching factor of the game and $d$ is
searched depth~\cite{ZHONG}.

The most time is spared if the pruning children of nodes are listed as first.
However in the worst case if children are sorted in opposite order to the
optimal, the whole minimax tree is searched. Consequently it is very important
to have nodes ordered well.


\section{Description of the Monte Carlo Tree search}
Monte Carlo methods were formerly used for approximation of mathematical,
physical, biological and others processes where full calculation would be
difficult or even impossible. For example in mathematics it is used for numeric
integration or estimate the value of~$\pi$. In games without perfect
information, such as poker, backgammon, or scrabble, using randomness has shown
to be beneficial too~\cite{MonteCarloMethod, MonteCarloGo}.

In game of Go was always a huge troublesome to build efficient evaluation
function. Unlike poker or backgammon, Go is a game with perfect information and
hence it is not so natural to use Monte Carlo methods. The first attempts to
use random approach as evaluating approximation of the Go position were in
1993. We present generalised Bernd's algorithm~\cite{BERND,KOZELEK}:

\begin{enumerate}
\item Play random game from given position to the end, with one exception,
	choose only steps not filling eyes. At the end of simulation count in the
	result of simulation for the first step played.
\item If there is time left go to 1.
\item Choose move with highest ratio between number times the move won when it
	  was played and the number of times it was played.
\end{enumerate}

However Bernd showed that after some time his algorithm indicates no further
improvement. It was necessary to realize that the main problem is that
algorithm searches all branches with the same probability. It should give more
attention to more promising branches. The question how to find fine balancing
between exploration and exploitation is classical mathematical Multi-armed
Bandit problem.

\subsection{Bandit Problem}
A $K$-armed bandit, is slot machine with $K$ arms. When arm is drawn it
produces reward. Distribution of each arm reward is independent on other arms
and previous draws of this arm. The task is to choose best strategy to maximize
sum of reward through iterative plays~\cite{MoGo,MultiarmedBandit}.

In~\cite{MultiarmedBandit} is presented three UCB1 algorithm (where UCB stands for Upper Confidence Bounds) for Bandit Problem:

\begin{enumerate}
\item Play each arm of the bandit once.
\item Play arm maximizing the formula $\overline X_i + \sqrt{2 \log n \over n_i}$,
	  where $\overline X_i$ is average value of the arm $i$, $n$ is number
	  of games that were played by parent of the $i$ and $n_i$ is number of
	  games played with arm $i$.
\end{enumerate}


\subsection{UCT algorithm}
UCT is abbrevation for  Upper Confidence bound to Trees. Shortly described it
is the UCB algorithm applied to minimax trees. The main idea is to consider
each node of minimax tree as multiarmed bandit problem and each child of the
node as independent arm~\cite{MoGo}.


\subsection{MCTS}
Monte Carlo Tree search is generic best-first-search algorithm applied to
trees which uses a random simulation as evaluation scheme. It consist of four
steps which are repeated as long as there is time left. The algorithm starts
with a tree containing only a root node representing starting
position~\cite{progressive-strategies,KOZELEK}.

The four mentioned steps are explained in the following list:

\begin{enumerate}
\item \emph{Selection:} Until a leaf node is reached a child of node is
recursively selected by using best-first manner. The node selection is usually
handled by UCT algorithm.
\item \emph{Expansion:} The selected leaf's children are generated.
\item \emph{Simulation} (also called \emph{playout}): A random game is played
from given position to certain depth.
\item \emph{Backpropagation:} The result of simulated games is stored into all
nodes visited during this iteration of algorithm in the Selection phase.
\end{enumerate}

Classical approach of the Simulation phase in Go bot programming is to generate
only steps not filling eyes and to stop playout when the end of game is
reached (which is not possible in Arimaa)~\cite{MoGo}.

The result of the Simulation phase is computed from the last position of the
simulation by a evaluation function. The probabilistic attitude of UCT/UCB1
requires the value of evaluation to be scaled in interval
$[-1,1]$~\cite{KOZELEK}.

A pseudocode of the full algorithm is shown in~\ref{alg:mcts} (inspired
by~\cite{MoGo}).

\lstset{language=Python, caption=Pseudocode of the MonteCarlo Tree Search, label=alg:mcts}
\begin{lstlisting}
playOneIteration():
    node[0] = rootNode
    i = 0
    while is_not_leaf(node[i]):
        node[i+1] = descendByUCB1(node[i])
        i = i + 1

    expandNode(node[i])
    node[i].value = getValueByMonteCarlo(node[i])

    for j in {0, ..., i-1}:
        updateValue(node[j],node[i].value)

while is_time():
    playOneIteration()

print descendByUCB1(rootNode)
\end{lstlisting}

We modified Kozelek's UCB exploration formula to form:
	$$
	\overline X_i + c \sqrt{\frac{log~n}{n_i}} + \frac{h_i}{n_i} + \frac{hh_i}{n_i}
	$$
where $\overline X_i + c \sqrt{\frac{log~n}{n_i}}$ is a generalised UCB1
formula, $h_i$ is heuristics evaluated for step leading to this position, and
$hh_i$ is history heuristic value.

One of the negatives of UCT algorithm is that UCB1 formula itself presumes for
the involved random variables to be identically distributed and independent,
which is not true in UCT~\cite{MoGo}. It was shown that the probability of
selecting the correct move quickly converges to 100\%~\cite{UCTanalysis}.


\chapter{Used optimisations in search engines}

In order to write a strong playing program in given game, it is necessary to
enhance chosen algorithm with various extensions. In this chapter we describe
the most used extensions for AlphaBeta or MCTS algorithm.

\section{AlphaBeta}
\subsection{Transposition Table}\label{AlphaBeta:TT}
If we look at the game tree of Arimaa, Chess or Go there is so many repetitions in
nodes of the tree for almost every position on the board. The Transposition table
is used to reduce the number of repetitions in tree.

COMPLETELY REWRITE: If in TT occurs result from previous search that is
shallower than we need, we prefer steps from Principal Variation of given entry
in succeeding search and use also its bounds.

If we find a satisfying result in TT we use its full PV and its score as ours.

\subsection{Iterative Deepening Framework}
Because there is no way how to determine how long the AlphaBeta search to
given depth will take, we need to find some time management tool. This is a
method in which we are iteratively starting new deeper and deeper searches.
Thanks to the exponential growth of the minimax tree with increasing depth we
know that a searching one level more will cost us approximately a lot more than
previous shallower searching.

Because we can sort moves using information gained from previous shallower
searches Iterative deepening used with other optimisation methods such as
History Heuristic or Transposition Table often cause a time save if we compare
it to just searching to a maximal possible depth~\cite{COX}.

\subsection{Aspiration Windows}
% TODO how to mention it?
% \subsection{Quiescent search}
%   Is nice way, how to reduce horizont effect ...
%   % [Related to arimaa: http://arimaa.com/arimaa/forum/cgi/YaBB.cgi?board=devTalk;action=display;num=1122418533] <- TODO: Maybe implement this way
%   % [http://mediocrechess.sourceforge.net/guides/quiescentsearch.html]

\subsection{Move ordering}
The following methods change the order in which branches are selected and then
inspected. It is very important for the AlphaBeta search to have nodes well
ordered. Because earlier we find a pruning child of a node the shorter time
we spend in it.

	\subsubsection{History heuristics}
	The main idea behind this extension is that if some step is good enough
	to cause so many pruning anywhere in the search tree and if it is valid in
	given position it could be also good here.

	It is implemented in way that we create a table with a score for every
	element from combination of players, piece, position and direction. During
	the AlphaBeta search we increase the score for element every time such
	combination causes pruning or finds the best score. It is believed that
	the deeper cut-off happens the more relevant it is and therefore the score
	is incremented by $d^2$ or $2^d$ where $d$ is a actual searched depth.

	During searching are nodes of the minimax tree sorted by score from the
	History Heuristic table in decreasing order~\cite{COX}.

	\subsubsection{Killer moves}
	When a step prunes branches in some position it is very natural to ask if
	the same step could cause pruning in another branch and the same depth of
	the tree. To go even further we take two last steps caused pruning to be
	preferred in the search. As Zhong found out three or more Killer moves would help~\cite{ZHONG}. TODO

	\subsubsection{Null move}
	TODO


\subsection{Heuristics}
All preceding optimisations do not change result given by the algorithm. They
could only significantly decrease amount of time needed to obtain such result.

The consequent optimisations are rather heuristics, because they try to
give you approximately good result in much shorter time. In this work, we
will not use those optimisations.

TODO Introduce  Negascout/PVS and MTD-f or delete this subsection.


\section{Monte Carlo Tree Search}
In AlphaBeta search we made great effort in sorting nodes of the searched tree
properly. In MCTS we need to use optimisations which helps us to gain better
and more informations from each iteration of algorithm on top of that.

We chose to implement only heuristic Tomáš Kozelek described in~\cite{KOZELEK}
as useful.

\subsection{Transposition table}
The motivation is the same as is in the AlphaBeta algorithm
(see~\ref{AlphaBeta:TT}). However in MCTS we are increasingly building game
tree instead of just exploring branches to some depth. A natural use of
Transposition table for MCTS is to share statistics for the same transpositions
in built game tree.

To do so we bind nodes considered the same to one. Bound nodes share their
children nodes, visit count and score statistic. We say that two nodes are the
same if they are in the same depth in minimax tree and if they represent the
same transposition. The game tree became Directed Acyclic Graph.

In implementation during the computation we keep table of all transpositions
and when any node is expanded we bind it to transposition in table if exists.

In Kozelek's work is regarded to be dangerous to bound visit count and score
statistic for children nodes~\cite{KOZELEK}. Nevertheless we believe that if
some step leads us to a position which is proven to be not worth trying in some
branch the same stands for all its transpositions.

\subsection{Progressive bias}
Progressive bias technique is a nice way how to combine offline learned
informations with online learned informations. The more we go through a node
the importance of offline learned information goes down and the importance of
online learned information became superior.

To do so ${H_B \over n_i}$ is added to the UCB formula. Where $H_B$ is
progressive bias coefficient computed by step-evaluation
function~\cite{progressive-strategies}.

In Arimaa such step-evaluation function should appreciate steps with Elephant
moving, killing an opponents piece, around previous steps, which are pushing or
pulling or making goal. Such function should also handicap players own piece
sacrifice steps or inverse steps~\cite{KOZELEK}.

\subsection{History heuristics}
Adds to UCB1 formula $+ {hh_i \over n_i}$. Where $hh_i$ is history
heuristics coefficient as is described in Kozeleks thesis~\cite{KOZELEK}.

TODO

Similar approach is All-Moves-As-First Heuristics however it updates statistics
only for all steps that were used in this simulation (in each node) (TODO)

\subsection{Best-of-$N$}
In random simulations, we may want to sacrifice true randomness for gaining
more objective results from playouts~\cite{HeavyPlayouts}. $N$ random steps are
generated and a step with the best value given by the same step-evaluation
function as was used in Progressive bias is chosen.

As a consequence, the number of playouts decreases, but the quality of
information learned in playouts improve and therefore the strength of the
program improve also.

\subsection{Children caching}
Kozelek introduced natural method how to decrease an amount of time spent in
node during selection part of the MCTS.

After some number of visits of a node its children caching is switched on. Then
a few best children are chosen and cached and every time the selection part of
the MCTS goes through this node it chooses a descend node from earlier cached
children. After some time it is necessary to discard and fill cache
again~\cite{KOZELEK}.

Using this optimisations should improve the speed of algorithm without negative
impact on quality.

\subsection{Maturity threshold}
Is another technique how to shorten time used in node selection. We
expand only nodes which had at least $threshold_maternity + depth\_of(node)$
visit count.

\subsection{Virtual visits}
In node expansion initialise new node with $v$ virtual visits. This small
change increases the power of the algorithm significantly. Kozelek
experimentally determined the best performance of algorithm for $v \in
[4,5]$~\cite{KOZELEK}.

\section{Independent optimisations}
The following extensions are considered must have in every successful Arimaa
bot and do not depend on used algorithm.

	\subsection{Zobrist keys}
	Motivation ...

	For every element of $Piece\times Player\times Position$ the random 64 bit
	number is generated. ...

	How nodes are stored in Transposition table entries

	\subsection{Bitboards}
	Because Arimaa could be played with standart chess set it is also possible
	to use from world of chess well known board representation called
	Bitboards. One may say that bitboards suits even better for Arimaa than for
	Chess.

% TODO
% \section{Comparison of used optimisations}
% 
% \subsection{History heuristic}
% In MCTS we prefer moves that were often chosen in going through tree. In
% AlphaBeta value of move is increasing when branch using selected step is
% pruned. Every time we order branches of the search tree by value given by
% history heuristics. This approach should increase amount of pruning.
% 
% \subsection{Transposition table}


\chapter{Implementation}
We developed set of libraries for play Arimaa to be used with both MCTS and
AlphaBeta algorithm using haskell as programming language. On top of those
libraries we built mentioned algorithms. The critical parts like bit operations
and evaluate function were written in C.

Haskell differs from other common used languages whith its lazynes. Which means
that only values that are truly needed for computation are evaluated. For
example we may define infinite list of prime numbers and then ask for third
element of this list which causes only the first three prime numbers to be
computed.

This paradigm is applied naturaly in the AlphaBeta search. When we are
exploring some node and in one of the first child's the cut-off occurs we
spared time not generating others.

% In a tree node during MCTS when children hashing mechanism is switched on,
% instead of alocating more space for $C$ cached nodes, actual nodes are lazily
% sorted in decreasing order and only $\mathcal{O}(C\log n)$ sorting realy
% hapends.

In Haskell it is possible to program in much higher level than it is in most
other programming languages. However it is also harder to reason about
performance. We are pretty sure that more experienced Haskell programmer would
write both engines more efficiently.

\subsubsection{Parameter tuning}

We made huge effort to keep our program rather simple and modular as much as
possible. Therefore one can switch on or off a lot of mentioned search
extension. In step generator we let the possibility to generate Pass step
switched off by default.

Aspiration window gave sometimes strange results so we left it unused by default.
Using history heuristic in Alpha Beta program tends to decrease quality of our
program and we believe that causing list of steps from given position to be
evaluated is limiting in comparison to have them evaluated lazily. However it
is known that importance of history heuristic grows when the depth increaseand
hence for mor efficient programs HH is much more important~\cite{ZHONG}.

Our AlphaBeta algorithm lack of standart Quiescence extension with trap control
or the Goal check. We believe that both the Trap control and the Goal check
could be somehow also included in MCTS.

\section{Parallelization}
In Monte Carlo Tree Search, make bot to be parallel is much more natural than
it is in AlphaBeta search. For each CPU one thread is created and there is only
one shared tree for all of them.

Haskell gave us easy threading and data structure locking capabilities.
Local mutexes are rather simple thanks to MVar structure.



\section{Evaluation function}
We wrote evaluation function completely in C to be as fast as possible.

In our work we do not focus on creating strong evaluation function even though
in order to have arimaa playing program, we need to have one. Therefore we
developed one not so strong and with possibility to easily replace it.

Our evaluation function lack of a very important part of standart evaluation
function -- the goal check.
Also it do not have frame and hostage detection.

\chapter{Methodology}

At the begging of the work we decided to follow the consequent schedule:

\begin{enumerate}% {0}
\item Study possible and most used method used in Arimaa, Chess and Go bot development.
\item Develop basic versions of bots with plain versions of AlphaBeta search
and Monte Carlo Tree Search. However a work on bots is never ending so we will
try to develop two comparable bots and try to measure how changes in settings
affect their win:loss rate.
\item Improve both programs with enhances learned in first part of the process and tune variables of those enhances in order to improve the quality of both programs as much as possible.
\item Compare those engines playing offline matches and analyse the meaning of
obtained results.
% \item Analise engines time cost centres. Where are bottlenecks.
\end{enumerate}

\noindent In order to obtain as much informations as possible we decided to run
tests consisting of these parameters:

\begin{enumerate}
\item Test bots and enhances they uses by playing games with 3, 10, or 30 seconds
time limits per turn to test how different time limits affects their efficiency.
\item Playing games with Transposition Tables' size limited to 100MB, 200MB, or
400MB.
\item Playing engines using one, two, four, or eight cores.
% \item run AB to same level of search and give MCTS the same amount of time to play
\end{enumerate}


\chapter{Results}
We ran our tests on machines 12 machines with
Intel\textregistered~Core\texttrademark~i7 CPU 920  @ 2.67GHz and 14 machines
with Intel\textregistered~Core\texttrademark2 Quad CPU Q9550 @ 2.83GHz. In all
general test cases, bots had 200MB of memory and time limit 3, 10 or 30 seconds
per turn. The logs from tests are saved on attached CD.

In each experiment we performed around 400 tests. Where the around means that
more than 420 tests started, but only result from tests ending in time were
counted. In all figures the solid curve named \texttt{full} represents the
number of games in percentages in which the MCTS using all extensions defeated
the AlphaBeta search bot also with all extensions enabled.

\begin{figure}[htb]
	\centering
	\includegraphics[width=0.9\textwidth]{graphs/mctsTuning.eps}
	\caption{MCTS optimisations' performance}
	\label{pic:mctsTuning}
\end{figure}

In Figure~\ref{pic:mctsTuning} is shown how disabling some MCTS's
optimisation affects the winning ratio between algorithms. The noHH stands for
switching off History heuristic in MCTS, and noHeavyPlayout stands for
switching off Best-of-N optimisation. It seems from the graph that we achieved
constant improvement just by switching History heuristics on. On the other
hand, the importance of using Heavy Playouts increases with bigger turn limit.

\begin{figure}[htb]
	\centering
	\includegraphics[width=0.9\textwidth]{graphs/generalTuning.eps}
	\caption{Comparing MCTS's and AlphaBeta's general capabilities}
	\label{pic:generalTuning}
\end{figure}

In Figure~\ref{pic:generalTuning} we can see how switching off Transposition
Tables (as \texttt{noTT} in graph) or changing evaluation function to function
used in bot\_Fairy (\texttt{EVAL=fairy} in graph) changed the winning ratio
between our two bots.

In this comparison, it is important to mention how using Fairy's evaluation
function affects bots strength. We performed another tests with 10 seconds per
turn time limit, in each either AlphaBeta or MCTS version of our bot uses
Fairy's evaluation function. Again the information about winning ratio is shown
from MCTS's point of view:

\begin{center}
\begin{tabular}{rl}
AlphaBeta with Fairy's evaluation function vs full MCTS: & 41.51\% \\
full AlphaBeta vs MCTS with Fairy's evaluation function: & 20.12\% \\
\end{tabular}
\end{center}

\noindent This means that using Fairy's evaluation is making our bots weaker.
Therefore we can conclude from the graph in Figure~\ref{pic:generalTuning} that
MCTS behaves better with worse evaluation functions compared to AlphaBeta.

From Figure~\ref{pic:generalTuning} we see that MCTS without Transposition
Tables defeats AlphaBeta more likely with bigger time limits. A possible
explanation is that either Transposition Tables are far more important in
AlphaBeta or using History Heuristic in MCTS compensates the loss of
Transposition Tables.

At the end of testing process we extended our evaluation function with
Kozelek's goal check from his bot Akimot. From Figure~\ref{pic:generalTuning}
we see that using goal check helped AlphaBeta a lot. Very likely, this
corresponds to the fact, that our AlphaBeta search in 30 seconds per turn
explore minimax tree to depth of eight steps more often.

For the next tests we changed the parameters. Both algorithms used all extensions
(without goal check), and time limit per turn was set to 15 seconds. 

\begin{figure}[h]
	\begin{center}$
	\begin{array}{cc}
	\includegraphics[width=2.7in]{graphs/cores.eps} &
	\includegraphics[width=2.7in]{graphs/memory.eps}
	\end{array}$
	\end{center}
\caption{Comparison using parallelization or different memory size}
\label{pic:cpuAndMemory}
\end{figure}

Figure~\ref{pic:cpuAndMemory} displays how the winning ratio
MCTS AlphaBeta search changed with increasing number of CPU's or amount of
memory. The CPU's graph represents how the ratio changed with one,
two, four or eight CPU's and the size of Transposition Tables fixed to 200~MB.
In the Memory size graph it is shown how the ratio
changed with a size limit of Transposition Tables set to 100~MB, 200~MB or
400~MB and the number of CPU's set to one.

For better explanation of the CPU's graph from Figure~\ref{pic:cpuAndMemory} we
looked how the number of iterations of MCTS algorithm and the depth of
AlphaBeta search in steps changed with added CPU's:

\begin{center}
\begin{tabular}{ccc}
Cores: & MCTS iterations: & Average depth of the AlphaBeta search in steps: \\
1      &  4566            & 6.9  \\
2      &  8370            & 8.3  \\
4      & 16161            & 8.8  \\
8      & 22083            & 8.8
\end{tabular}
\end{center}

Possible explanations of graph in~\ref{pic:cpuAndMemory} could be following:
\begin{itemize}
\item The results from CPU's graph could be  disadvantageous for the AlphaBeta search due to some artifact of the implementation.
\item AlphaBeta uses mutexes to lock global Transposition Tables during
insertions. MCTS uses mutexes for nodes, History heuristic table and
Transposition Tables locking. However MCTS spends most of the time in playouts
by generating and evaluating steps so it is not affected by the lockings so
much (See Appendix~\ref{statistics} for detailed statistics).
\item With limited memory resources (the 200~MB memory limit for Transposition
Tables in our case) having more CPU's becomes superior.
\item Light static evaluate function helps a lot with increasing number of CPU's.
\item For the AlphaBeta version of our bot, having more than 200~MB for
Transposition Tables with 15 seconds per turn time limit does not improve the
quality of the bot so much.
\item The previous conclusion about AlphaBeta's memory needs corresponds with these results.
\end{itemize}

% \chapter*{Conclusion}
% \addcontentsline{toc}{chapter}{Conclusion}

\chapter{Results}
Writing two programs at once turns out to be more difficult than we thought. In
our development we have not focussed on high time optimisations and
sophisticated goal check and trap control in position evaluation. These three
thinks are considered the most difficult in Arimaa bot development.

\section{Conclusion}
Using MCTS as is or as was provided by Kozelek is not enought for bot
programming in Arimaa, but UCB formula or UCT algorithm itself can be valuable
in yet unknown hybrid algorithm using both AlphaBeta and MCTS algorithms.

??? Combine previous and consequent paragraphs ???

The most prommising variant of the Arimaa playing program using MCTS algorithm
could be some combination of MCTS and AlphaBeta. Using full four depth
AlphaBeta search while expanding leafs of UCT. (after node expansion make
tactiacal lookahead)

REWRITE: We also believe that using more sophisticated board representation
would help MCTS algorithm against AlphaBeta, because the cost of managing
clever representation returns with more times you use the clever
representation, which dominates in MCTS. (From one position MCTS often
generates possible moves more than once in contrast of AlphaBeta).

\section{Further work}
% TODO it should be further work from our point of view (to attack another similar topics to ours)

\begin{enumerate}
\item Find way how to implement some kind of Quiescence search to MCTS and compare with similar approach in AlphaBeta.
\item ...
\item Optimise playouts using more sophisticated way how to generate steps. Interesting ideas are described in Zhong's work.\cite{ZHONG}
\end{enumerate}

The quality of bots can be improved by:
\begin{enumerate}
\item create better evaluation function,
\item try progressive pruning methods\cite{progressive-strategies},
\item use more optimised data structures and random number generator,

Goal check??

\end{enumerate}


\appendix
\def\bibname{Literature}
\begin{thebibliography}{99}
\addcontentsline{toc}{chapter}{\bibname}

\bibitem{KOZELEK}
	{\sc KOZELEK,} Tomáš.
	\emph{MCTS methods in the game of Arimaa}.
	Master's thesis.
	Charles University in Prague, 2009.

\bibitem{COX}
	{\sc COX,} Christ-Jan.
	\emph{Analysis and Implementation of the Game Arimaa}.
	Master Thesis.
	Universiteit Maastricht, 2006.

\bibitem{arimaa.com}
	Arimaa homepage.
	\texttt{http://arimaa.com/}.

\bibitem{syed}
	Omar {\sc SYED} and Aamir {\sc SYED}.
	\emph{Arimaa: A New Game Designed to be Difficult for Computers}.
	Journal of the International Computer Games Association, 26(2):138–139, 2003.

\bibitem{knuth:alphabeta}
	Donald E. {\sc KNUTH} and R. W. {\sc MOORE} (1975).
	\emph{An Analysis of Alpha-Beta Pruning}.
	Artificial Intelligence Vol. 6, No. 4: 293-326. Reprinted in Selected Papers on Analysis of Algorithms by Donald E. Knuth, Published 2000, by Addison-Wesley
	ISBN: 1-57585-211-5.

\bibitem{ZHONG}
	{\sc Zhong,} Haizhi.
	\emph{Building a Strong Arimaa-playing Program}.
	Master Thesis.
	University of Alberta, 2005.

\bibitem{progressive-strategies}
	Guillaume {\sc CHASLOT}, Mark {\sc WINANDS}, Jaap H. van den {\sc HERIK},
	Jos {\sc UITERWIJK}, and Bruno {\sc BOUZY}.
	\emph{Progressive strategies for Monte-Carlo tree search}.
	In Joint Conference on Information Sciences, Salt Lake City 2007, Heuristic
	Search and Computer Game Playing Session, 2007.


\end{thebibliography}

[2] Peter Auer, Nicol` Cesa-Bianchi, and Paul Fischer. Finite-time analysis of
    the o multiarmed bandit problem. Mach. Learn., 47(2-3):235–256, 2002.
[BERND] Bernd Brugmann. Monte Carlo go. Technical report, 1993.
[15] Chaslot, Winonds, ..., Progressive strategies for Monte-Carlo Tree Search, 2007.


\chapwithtoc{List of Abbreviations/Glossary}
\begin{acronym}
	\acro{MCTS}{Monte Carlo Tree Search}
	\acro{UCB}{Upper Confidence Bounds}
	\acro{UCT}{Upper Confidence bound to Trees or in other words UCB applied to minimax trees}
	\acro{filling eye}{}
\end{acronym}


\chapter{Appendix 1}
TODO Memory/Heap and time/cost statistics.


\chapter{Appendix 2 - User documentation}
Source code is available in http://github.com/JackeLee/rabbocop or in included CD.

\section{Compiling options}
Additional make parameters:

VERBOSE=$n$
  For IterativeAB program, it will print actual best move and score anytime
  it searches to desired depth.
  For MCTS it will print actual best move and score after n iterations of
  UCT algorithm.

EVAL=fairy

JUDY=1
HASKELL\_HASH=1

TODO: NULL\_MOVE CAN\_pass

WINDOW=?
	It switches on aspiration window option for AlphaBeta algorithm.

noHH=1
	It disables history heuristics optimisations for MCTS.

abHH=1
	Enables history heuristics optimisations for AlphaBeta.

noHeavyPlayout=1
	This disables heuristic in playouts.

CORES=$n$ - TODO
	Set the number of available cores in computer. It is also needed to run
	binary with +RTS -N$n$ parameter.

PROF=1
	It enables profiling options.

runtest
	Can have

play


\openright
\end{document}
